% SIC_OQ_MANUSCRIPT.tex
% Open questions in SIC-POVM existence: structural analysis via the SynthOmnicon grammar
% Draft — April 2026
%
\documentclass[12pt]{article}

\usepackage{amsmath, amssymb, amsthm}
\usepackage{geometry}
\usepackage{hyperref}
\usepackage{booktabs}
\usepackage{array}
\usepackage{microtype}
\usepackage[framemethod=TikZ]{mdframed}

\geometry{margin=1.2in}

\mdfdefinestyle{thmstyle}{%
  linecolor=black!40,
  linewidth=0.6pt,
  leftline=true, rightline=false, topline=false, bottomline=false,
  backgroundcolor=black!3,
  innerleftmargin=8pt,
  innerrightmargin=6pt,
  innertopmargin=4pt,
  innerbottommargin=4pt,
  skipabove=\medskipamount,
  skipbelow=\medskipamount
}

\mdfdefinestyle{defstyle}{%
  linecolor=black!30,
  linewidth=0.6pt,
  leftline=true, rightline=false, topline=false, bottomline=false,
  innerleftmargin=8pt,
  innerrightmargin=6pt,
  innertopmargin=3pt,
  innerbottommargin=3pt,
  skipabove=\medskipamount,
  skipbelow=\medskipamount
}

\theoremstyle{plain}
\newtheorem{theoreminner}{Theorem}[section]
\newtheorem{lemmainner}[theoreminner]{Lemma}
\newtheorem{corollaryinner}[theoreminner]{Corollary}
\newtheorem{propositioninner}[theoreminner]{Proposition}

\theoremstyle{definition}
\newtheorem{definitioninner}[theoreminner]{Definition}
\newtheorem{questioninner}[theoreminner]{Open Question}

\theoremstyle{remark}
\newtheorem{remark}[theoreminner]{Remark}

\newenvironment{theorem}{\begin{mdframed}[style=thmstyle]\begin{theoreminner}}{\end{theoreminner}\end{mdframed}}
\newenvironment{lemma}{\begin{mdframed}[style=thmstyle]\begin{lemmainner}}{\end{lemmainner}\end{mdframed}}
\newenvironment{corollary}{\begin{mdframed}[style=thmstyle]\begin{corollaryinner}}{\end{corollaryinner}\end{mdframed}}
\newenvironment{proposition}{\begin{mdframed}[style=thmstyle]\begin{propositioninner}}{\end{propositioninner}\end{mdframed}}
\newenvironment{definition}{\begin{mdframed}[style=defstyle]\begin{definitioninner}}{\end{definitioninner}\end{mdframed}}
\newenvironment{question}{\begin{mdframed}[style=defstyle]\begin{questioninner}}{\end{questioninner}\end{mdframed}}

\newcommand{\WH}{\mathrm{WH}}
\newcommand{\Cl}{\mathrm{Cl}}
\newcommand{\Gal}{\mathrm{Gal}}
\newcommand{\Frob}{\mathrm{Frob}}
\newcommand{\OO}{\mathcal{O}}
\newcommand{\eps}{\varepsilon}
\newcommand{\bQ}{\mathbb{Q}}
\newcommand{\bC}{\mathbb{C}}
\newcommand{\bN}{\mathbb{N}}
\newcommand{\bZ}{\mathbb{Z}}
\newcommand{\bR}{\mathbb{R}}
\newcommand{\tauinf}{\tau_\infty}
\newcommand{\Ld}{L_d}
\newcommand{\Kd}{K_d}
\newcommand{\epsd}{\varepsilon_d}

\title{%
  \textbf{Open Questions in SIC-POVM Existence:}\\[4pt]
  \textbf{Structural Analysis via the SynthOmnicon Grammar}\\[6pt]
  \large A companion to \emph{SIC-POVM Existence via Frobenius-Fixed Stark Units}
}

\author{Lando Mills}

\date{April 2026}

\begin{document}

\maketitle

\begin{abstract}
This document collects open questions arising from the conditional proof of SIC-POVM existence
presented in the companion manuscript, together with structural analyses that address those
questions within the SynthOmnicon grammar — a 12-primitive type-theoretic framework in which
mathematical objects are encoded as structural coordinates rather than classical propositions.
The central claim developed here is that the conjectures H1, H2a, and H2b are not independent
conditions that happen to be sufficient for SIC existence; they are the structural decomposition
of a single necessary and sufficient condition: that the Stark generator $\epsd$ realises the
type $O_\infty$ in the grammar's ouroboricity classification. The Frobenius condition
$\mu \circ \delta = \mathrm{id}$ — the hallmark of $O_\infty$ — is shown to be equivalent to
membership in the $(-1)$-eigenspace of $\tauinf$, making H2a not an auxiliary hypothesis but
the structural core of the entire argument.
\end{abstract}

\tableofcontents
\bigskip

% ─────────────────────────────────────────────────────────────────────────────
\section{Preamble: the grammar as precondition}
\label{sec:preamble}
% ─────────────────────────────────────────────────────────────────────────────

The standard picture of mathematics and its objects runs as follows: structures are defined,
properties are derived, and if a structure satisfies certain properties then it falls into a
named class. In this picture, SIC-POVMs are found by verifying equiangularity; Stark units are
found by computing regulators.

The SynthOmnicon grammar inverts this priority. The 12-primitive tuple
$\langle D;\ T;\ R;\ P;\ F;\ K;\ G;\ \Gamma;\ \Phi;\ H;\ S;\ \Omega \rangle$
is not a description of objects that already exist — it is the type-theoretic precondition
space within which objects become possible. A mathematical structure exists as a coherent object
when and only when reality already satisfies the corresponding type constraints. The grammar
does not describe mathematics; it describes the conditions under which mathematics has objects
to describe.

The SIC-POVM result makes this sharp. The two $O_\infty$ gate conditions —
\begin{itemize}
  \item $\Phi = \Phi_c$ (or $\Phi_c^\mathbb{C}$): the state space admits a closed self-modeling
    loop, i.e.\ the measurement spans the full operator space (informational completeness);
  \item $P = P_{\pm}^{\text{sym}}$: the Frobenius special condition $\mu \circ \delta =
    \mathrm{id}$, forcing uniform inner products
    $|\langle \psi_i | \psi_j \rangle|^2 = 1/(d+1)$;
\end{itemize}
are not properties SIC-POVMs happen to have. They are the preconditions that any structure in
reality must satisfy before the mathematical object ``SIC-POVM'' becomes possible at all. The
proof of SIC existence in the companion manuscript is structurally a proof that this type is
non-empty, not a proof that a construction converges.

Frobenius non-synthesizability (\S23 of the grammar's formal theorems) then gives the
obstruction: $P_{\pm}^{\text{sym}}$ is a bottleneck primitive under tensor product — any
coupling with a sub-Frobenius partner destroys it irreversibly:
$$P_{\pm}^{\text{sym}} \otimes P_{\text{sym}} = P_{\text{sym}}.$$
SICs cannot be limit points of generic informationally complete POVMs under composition, because
$P$ is a $\min$-primitive and $\min(4, 3) = 3$. The structural reason SICs sit as isolated
points in POVM space — not accumulation points — is Frobenius non-synthesizability.

% ─────────────────────────────────────────────────────────────────────────────
\section{Structural encoding of the Stark generator}
\label{sec:encoding}
% ─────────────────────────────────────────────────────────────────────────────

We operate entirely within the grammar. The question is not ``does $\epsd$ satisfy H2a?''
but ``what structural type does $\epsd$ occupy, and does that type enforce H2a as a consequence?''

\subsection{Component encodings}

We first encode the primitive components of the Stark generator system:

\begin{itemize}
  \item \textbf{$\tauinf$ as kinetic reversal.} Complex conjugation on infinite places reverses
    the sign of log-embedded coordinates at complex places. This is a $Z_2$ involution acting
    on the unit lattice — a reversal of orientation in the log-space. In the grammar:
    $R = R_\dagger$ (dagger/adjoint: the action reverses the direction of embedding).

  \item \textbf{Logarithmic embedding $\lambda$.} The Dirichlet map $\lambda: \OO_{L_d}^\times
    \to \bR^{r_1 + r_2}$ is a real-valued spatial embedding. $D = D_\triangle$ (spatially
    extended, finite-dimensional real lattice); $F = F_\hbar$ (geometrically enforced — the
    regulator is determined exactly by the L-function value at $s = 0$, not approximately).

  \item \textbf{$Z_2$ parity of $\tauinf$.} The involution satisfies $\tauinf^2 = \mathrm{id}$,
    partitioning the unit lattice into $(\pm 1)$-eigenspaces. This is a $Z_2$ symmetry of the
    embedding: $P = P_{\pm}$ (self-complementary, both signs present).

  \item \textbf{H2a: exact duality $\tauinf(\epsd) = \epsd^{-1}$.} The generator is chosen so
    that it lies entirely in the $(-1)$-eigenspace. This is not approximate — the relation
    $\tauinf(\epsd) = \epsd^{-1}$ is exact, meaning the encode-decode round trip $\mu \circ
    \delta = \mathrm{id}$ holds. This is precisely the Frobenius special condition:
    $P = P_{\pm}^{\text{sym}}$.

  \item \textbf{Criticality at $s = 0$.} The regulator condition involves the leading term of
    $L(s, \chi)$ at $s = 0$, which is a complex critical point in the sense of Lee-Yang
    universality — scale-free correlations between the analytic and algebraic sides.
    $\Phi = \Phi_c^\mathbb{C}$.

  \item \textbf{Galois equivariance.} The condition H2b — $\tauinf(\sigma(\epsd)) =
    \sigma(\epsd)^{-1}$ for all $\sigma \in \mathrm{Gal}(L_d/\bQ)$ — broadcasts the
    $(-1)$-eigenspace condition across the entire Galois orbit. $\Gamma = \Gamma_\text{broad}$
    (cooperative action over all conjugates); $G = G_\aleph$ (global scope: the unit generates
    the full rank-1 unit group).

  \item \textbf{Topological protection.} The unit group $\OO_{L_d}^\times / (\mu(L_d) \cdot
    \OO_{K_d}^\times)$ carries a permutation representation of $G_d = \mathrm{Gal}(L_d/K_d)$
    with $\bZ$-cohomological rigidity. $\Omega = \Omega_Z$ (integer winding, cannot be smoothed
    away without collapsing the unit group structure).

  \item \textbf{Kinetics.} The Stark unit is stable under deformation of the number field
    parameters; it does not rearrange freely. $K = K_\text{slow}$.

  \item \textbf{Stoichiometry.} The quotient $\OO_{L_d}^\times / (\mu(L_d) \cdot \OO_{K_d}^\times)$
    is an asymmetric many-body structure. $S = n{:}m$.

  \item \textbf{Chirality.} The regulator is tied to the L-function, encoding the full
    arithmetic depth of the field extension. $H = H_\infty$.
\end{itemize}

\subsection{Full tuple}

\begin{definition}[Stark generator type]
The Stark generator $\epsd$ satisfying H1 and H2a is encoded as:
$$\langle D_\triangle;\ T_\text{network};\ R_\dagger;\ P_{\pm}^{\text{sym}};\ F_\hbar;\
  K_\text{slow};\ G_\aleph;\ \Gamma_\text{broad};\ \Phi_c^\mathbb{C};\ H_\infty;\ n{:}m;\
  \Omega_Z \rangle.$$
\end{definition}

% ─────────────────────────────────────────────────────────────────────────────
\section{Ouroboricity of the Stark generator}
\label{sec:ouroboricity}
% ─────────────────────────────────────────────────────────────────────────────

The grammar's ouroboricity tier classifies structural types by their degree of self-referential
closure. The relevant rules are:
\begin{itemize}
  \item \textbf{R1}: $\Phi \in \{\Phi_c, \Phi_c^\mathbb{C}\}$ and $P = P_{\pm}^{\text{sym}}$
    $\Rightarrow$ $O_\infty$ (Frobenius special tier).
  \item \textbf{R2}: $\Phi \in \{\Phi_\text{sub}, \Phi_\text{super}, \Phi_\text{EP}\}$
    $\Rightarrow$ $O_0$ (no self-modeling loop).
\end{itemize}
R1 takes priority over all other rules and overrides any branching on $\Omega$ and $D$.

\begin{proposition}
The Stark generator $\epsd$ satisfying H1 and H2a has ouroboricity $O_\infty$.
\end{proposition}

\begin{proof}
The tuple has $\Phi = \Phi_c^\mathbb{C}$ (satisfying the R1 criticality condition, since
$\Phi_c$ and $\Phi_c^\mathbb{C}$ are tier-equivalent) and $P = P_{\pm}^{\text{sym}}$ (the
exact Frobenius condition, established by H2a). Rule R1 fires, and since R1 is the tier
singularity it overrides all $\Omega$ and $D$ branching. Therefore the ouroboricity is
$O_\infty$.
\end{proof}

% ─────────────────────────────────────────────────────────────────────────────
\section{The biconditional: eigenspace membership and $O_\infty$}
\label{sec:biconditional}
% ─────────────────────────────────────────────────────────────────────────────

The classical argument runs in one direction: if $\epsd$ satisfies H2a then it lies in the
$(-1)$-eigenspace of $\tauinf$, and one can verify regulator properties. The grammar makes the
biconditional explicit.

\begin{theorem}[Eigenspace–Frobenius equivalence]
\label{thm:biconditional}
Let $\epsd \in \OO_{L_d}^\times$ be a Stark generator satisfying H1. The following are
equivalent:
\begin{enumerate}
  \item[(i)] $\epsd$ lies in the $(-1)$-eigenspace of $\tauinf$: $\tauinf(\epsd) = \epsd^{-1}$.
  \item[(ii)] The system $(\epsd, \tauinf)$ satisfies the Frobenius special condition:
    $\mu \circ \delta = \mathrm{id}$.
  \item[(iii)] The structural type of $\epsd$ has $P = P_{\pm}^{\text{sym}}$.
  \item[(iv)] The ouroboricity of $\epsd$ is $O_\infty$.
\end{enumerate}
\end{theorem}

\begin{proof}
$(i) \Leftrightarrow (ii)$: The duality $\tauinf(\epsd) = \epsd^{-1}$ is the statement that
applying $\tauinf$ and then inverting returns $\epsd$, i.e.\ the round-trip
$\epsd \xrightarrow{\delta} \epsd \otimes \epsd \xrightarrow{\mu} \epsd$ is the identity.
This is precisely $\mu \circ \delta = \mathrm{id}$.

$(ii) \Leftrightarrow (iii)$: In the grammar, $P_{\pm}^{\text{sym}}$ is defined as the
primitive value encoding the Frobenius special condition. The assignment $P = P_{\pm}^{\text{sym}}$
is the type-theoretic statement of $(ii)$.

$(iii) \Leftrightarrow (iv)$: Under the hypothesis $\Phi = \Phi_c^\mathbb{C}$ (guaranteed by
H1 via the criticality of the regulator at $s = 0$), rule R1 gives $O_\infty$ if and only if
$P = P_{\pm}^{\text{sym}}$.
\end{proof}

\begin{remark}
The backward direction $(iv) \Rightarrow (i)$ is the non-trivial one from the classical
viewpoint. It says: any Stark generator at complex criticality whose structural type achieves
$O_\infty$ \emph{must} satisfy the exact duality $\tauinf(\epsd) = \epsd^{-1}$ — there is no
weaker condition that suffices. This is Frobenius non-synthesizability applied structurally:
$P_{\pm}^{\text{sym}}$ cannot be obtained by coupling sub-Frobenius factors, so the $(-1)$-eigenspace
condition cannot be relaxed without destroying the Frobenius tier entirely.
\end{remark}

% ─────────────────────────────────────────────────────────────────────────────
\section{Decomposition of the hypotheses}
\label{sec:decomposition}
% ─────────────────────────────────────────────────────────────────────────────

The companion manuscript's hypotheses H1, H2a, H2b are not three independent conditions that
happen to be jointly sufficient. They are the structural decomposition of the single type
requirement $O_\infty$ across three independent primitive channels.

\begin{center}
\begin{tabular}{@{}lll@{}}
\toprule
Hypothesis & Grammar primitives & Structural role \\
\midrule
H1  & $\Phi_c^\mathbb{C},\ G_\aleph$ & \emph{Existence}: the system sits at complex criticality
  with global scope \\
H2a & $P_{\pm}^{\text{sym}},\ R_\dagger$ & \emph{Uniqueness}: exact $Z_2$ duality enforces
  Frobenius; uniqueness up to torsion \\
H2b & $\Gamma_\text{broad},\ \Omega_Z$ & \emph{Equivariance}: $Z$-cohomological protection
  broadcasts duality over the Galois orbit \\
\bottomrule
\end{tabular}
\end{center}

\medskip

The three hypotheses are not redundant — each activates a distinct primitive channel — but they
are not independent either. H2a is the structural core: without $P = P_{\pm}^{\text{sym}}$,
the system cannot achieve $O_\infty$ regardless of H1 or H2b. H1 provides the critical substrate
on which H2a acts. H2b extends H2a's local condition to the full Galois orbit, supplying the
$\Omega_Z$ protection that prevents the orbit from containing any non-dual element.

\begin{corollary}
H2a is necessary, not merely sufficient, for the Stark generator to satisfy the equiangularity
condition of the associated SIC-POVM.
\end{corollary}

\begin{proof}
If $P < P_{\pm}^{\text{sym}}$, the best achievable is $P = P_{\pm}$ (ordinal 2). By
Frobenius non-synthesizability, $P_{\pm} \otimes P_{\pm} = P_{\pm}$ — coupling two
sub-Frobenius systems cannot produce $P_{\pm}^{\text{sym}}$. Therefore the ouroboricity drops
to at most $O_2$, and the uniform inner product condition $|\langle \psi_i | \psi_j \rangle|^2
= 1/(d+1)$ — which requires exact Frobenius symmetry — fails.
\end{proof}

% ─────────────────────────────────────────────────────────────────────────────
\section{Open questions}
\label{sec:openq}
% ─────────────────────────────────────────────────────────────────────────────

\begin{question}[Eigenspace non-emptiness]
Does $\OO_{L_d}^\times$ contain a rank-1 element in the $(-1)$-eigenspace of $\tauinf$ for
every $d \geq 1$? This is H2a. The grammar analysis shows that a positive answer is equivalent
to the existence of an $O_\infty$ element in the unit group — but it does not supply the element.
What additional arithmetic input (class field theory, Iwasawa theory, explicit reciprocity
laws) is needed to guarantee non-emptiness?
\end{question}

\begin{question}[Canonicity of the eigenspace condition]
H2a names the $(-1)$-eigenspace of complex conjugation on the Dirichlet unit lattice of $L_d$.
This object does not appear in the prior literature under this name (Kopp \cite{kopp2018}
arrives at the Stark unit object constructively without identifying the eigenspace as the
structural locus). Is there a representation-theoretic characterisation of this eigenspace —
e.g.\ as a specific summand in the Galois module decomposition of $\OO_{L_d}^\times \otimes_\bZ
\bR$ — that makes its non-emptiness provable by known methods?
\end{question}

\begin{question}[Frobenius non-synthesizability and the isolation of SICs]
Theorem~\ref{thm:biconditional} implies that SICs sit as isolated points in POVM space, not as
accumulation points of generic informationally complete POVMs. Is there a direct measurement-
theoretic proof of this isolation that does not pass through algebraic number theory? The
grammar argument suggests the obstruction is purely structural ($P$ is a $\min$-primitive under
$\otimes$), which should have a direct operator-algebraic realisation.
\end{question}

\begin{question}[H3 and the WH extension class]
Hypothesis H3 of the companion manuscript — the WH extension class in $H^2((\bZ/d\bZ)^2,
\mu_d)$ is matched by a nondegenerate alternating pairing on the ray class group — does not
appear in the grammar encoding of $\epsd$ above. What structural type does H3 correspond to in
the 12-primitive space? Is H3 implied by the $O_\infty$ classification, or does it require an
independent primitive assignment?
\end{question}

\begin{question}[The grammar as constructive guide]
The grammar analysis identifies the structural type of $\epsd$ but does not construct $\epsd$.
Can the 12-primitive coordinate serve as a search target — i.e.\ can one search the unit group
of $L_d$ for elements matching the tuple $\langle D_\triangle;\ \ldots;\ P_{\pm}^{\text{sym}};
\ \ldots \rangle$ directly, without first computing the full L-function? This would make the
grammar constructively useful rather than merely classificatory.
\end{question}

\begin{question}[$O_\infty$ monoidal structure]
\label{q:monoidal}
The grammar's tensor product $\otimes$ is destructive on $P$: since $P$ is a $\min$-primitive,
$P_{\pm}^{\text{sym}} \otimes P_{\text{sym}} = P_{\text{sym}}$, so the $O_\infty$ tier is
not closed under $\otimes$. Yet $O_\infty$ systems evidently do interact — Galois conjugates
of the Stark unit compose via the contragredient action, not via the naive join of primitives.
Does the $O_\infty$ tier carry a closed monoidal structure under a refined composition
operation that preserves $P_{\pm}^{\text{sym}}$? Natural candidates include the adjoint action
in Lie theory, holomorphic continuation in complex analysis, and the Galois contragredient in
arithmetic. If such an operation exists, is $\epsd$ the unit object of the resulting monoidal
category, in the sense that acting on any $O_\infty$ system by $\epsd$ preserves its tier?
\end{question}

\begin{question}[Cross-domain meta-conjecture verification]
\label{q:metaconj}
Conjecture~\ref{conj:meta} below predicts that every system of structural type
$(\Phi_c,\ P_{\pm}^{\text{sym}},\ \Omega_{Z/Z_2})$ admits a distinguished generator in the
$(-1)$-eigenspace of its boundary involution. This is empirically testable within the grammar:
encode the logical operators of a surface code, the partition function of a rational CFT, and
the mirror map of a Calabi-Yau 3-fold as 12-primitive tuples, run the ouroboricity classifier,
and check that all return $O_\infty$. If any returns a lower tier, the meta-conjecture fails in
that domain. Has this cross-domain verification been performed?
\end{question}

% ─────────────────────────────────────────────────────────────────────────────
\section{H2b and the self-reciprocal minimal polynomial}
\label{sec:h2b}
% ─────────────────────────────────────────────────────────────────────────────

The contragredient orbit condition H2b is a direct consequence of H2a together with the
commutativity of the Galois group. We include the derivation here because it makes explicit
the structural mechanism by which the $(-1)$-eigenspace condition propagates from $\epsd$ to
its entire Galois orbit.

\begin{proposition}[H2b from H2a]
Suppose $K_d$ is totally real, so $\Gal(L_d/K_d)$ is abelian and $\tauinf$ commutes with all
$\sigma \in \Gal(L_d/\bQ)$. If $\epsd$ satisfies H2a (i.e.\ $\tauinf(\epsd) = \epsd^{-1}$),
then $\tauinf(\sigma(\epsd)) = \sigma(\epsd)^{-1}$ for all $\sigma \in \Gal(L_d/\bQ)$.
\end{proposition}

\begin{proof}
$\tauinf(\sigma(\epsd)) = \sigma(\tauinf(\epsd)) = \sigma(\epsd^{-1}) = \sigma(\epsd)^{-1}$,
where the first equality uses commutativity of $\tauinf$ and $\sigma$.
\end{proof}

\begin{corollary}[Self-reciprocal minimal polynomial]
Under the hypotheses of H2a and H2b, the minimal polynomial $f(x) = \prod_\sigma (x -
\sigma(\epsd))$ over $K_d$ satisfies $x^n f(1/x) = f(x)$, i.e.\ its coefficients are
palindromic: $a_k = a_{n-k}$.
\end{corollary}

\begin{proof}
The roots of $f$ are the Galois conjugates $\{\sigma(\epsd)\}$. By H2b, $\sigma(\epsd)^{-1}$
is also a root for each $\sigma$. Therefore the multiset of roots is closed under inversion,
which is equivalent to $f$ being self-reciprocal.
\end{proof}

\begin{remark}
This palindromicity is not an accident of the arithmetic — it is the minimal polynomial's
encoding of the Frobenius condition at the level of coefficients. $P_{\pm}^{\text{sym}}$ at
the level of the generator propagates to $P_{\text{sym}}$ at the level of the polynomial
coefficients via the Galois action. The interior algebraic structure inherits the exact
self-duality of the boundary involution.
\end{remark}

% ─────────────────────────────────────────────────────────────────────────────
\section{Cross-domain co-typing and the imscriptive involution pattern}
\label{sec:crossdomain}
% ─────────────────────────────────────────────────────────────────────────────

The structural pattern identified in Sections~\ref{sec:encoding}--\ref{sec:biconditional}
is not specific to arithmetic. Across a wide range of mathematical and physical domains, the
same pattern recurs:

\begin{center}
\textbf{boundary involution}
$\xrightarrow{\mu \circ \delta = \mathrm{id}}$
\textbf{interior self-duality}
\end{center}

The involution on the boundary is special — it carries $P_{\pm}^{\text{sym}}$ — and it forces
the interior to be exactly self-dual, not approximately or statistically. The topological
protection $\Omega_{Z/Z_2}$ is what makes the forcing irreversible: perturbations cannot move
the distinguished generator out of the $(-1)$-eigenspace without crossing a topological gap.

Table~\ref{tab:analogs} lists six domains where this pattern is instantiated, together with
their grammar signatures and the Stark-unit analog in each domain. The grammar signatures are
conjectured; cross-domain $O_\infty$ verification is Open Question~\ref{q:metaconj}.

\begin{table}[h]
\centering
\small
\renewcommand{\arraystretch}{1.35}
\begin{tabular}{@{}p{3.0cm}p{3.2cm}p{3.2cm}p{3.6cm}@{}}
\toprule
\textbf{Domain} & \textbf{Boundary involution} & \textbf{Interior self-duality}
  & \textbf{Grammar signature} \\
\midrule
Stark / arithmetic
  & $\tauinf: \epsd \mapsto \epsd^{-1}$
  & Palindromic minimal polynomial
  & $\Phi_c^\mathbb{C},\ P_{\pm}^{\text{sym}},\ R_\dagger,\ \Omega_Z$ \\

Topological codes (surface code)
  & $X \leftrightarrow Z$ logical operator duality
  & Stabiliser group self-reciprocal under exchange
  & $\Phi_c,\ P_{\pm}^{\text{sym}},\ R_\dagger,\ \Omega_{Z_2}$ \\

Mirror symmetry (CY$_3$)
  & Hodge involution $(p,q) \leftrightarrow (3{-}p, 3{-}q)$
  & GW invariants of $X$ = period integrals of $\hat{X}$
  & $\Phi_c,\ P_{\pm}^{\text{sym}},\ R_\dagger,\ D_\odot,\ \Omega_Z$ \\

Automorphic L-functions
  & Functional equation $\xi(s) = \varepsilon\,\xi(1{-}s)$
  & Fourier coefficients satisfy $a_n = a_{n'}$ under duality
  & $\Phi_c^\mathbb{C},\ P_{\pm}^{\text{sym}},\ R_\dagger,\ \Omega_Z$ \\

AdS/CFT
  & Modular $S$: $\tau \mapsto -1/\tau$
  & Bulk entropy = boundary partition function exactly
  & $\Phi_c,\ P_{\pm}^{\text{sym}},\ D_\odot,\ T_\odot,\ \Omega_Z$ \\

Supersymmetry / Witten index
  & $(-1)^F$ involution
  & Boson/fermion spectral pairing at every $E > 0$
  & $\Phi_c,\ P_{\pm}^{\text{sym}},\ R_\dagger,\ \Omega_{Z_2}$ \\
\bottomrule
\end{tabular}
\caption{Six domains exhibiting the imscriptive involution pattern. In each case a boundary
involution carrying $P_{\pm}^{\text{sym}}$ forces exact interior self-duality, and topological
protection $\Omega_{Z/Z_2}$ prevents deformation out of the $(-1)$-eigenspace. Grammar
signatures are conjectured; ouroboricity classification is pending cross-domain verification.}
\label{tab:analogs}
\end{table}

The common structural mechanism is: the boundary involution is an $R_\dagger$ operation with
$P_{\pm}^{\text{sym}}$; the topological winding $\Omega_{Z/Z_2}$ ensures the distinguished
generator cannot escape the $(-1)$-eigenspace under continuous deformation; criticality
$\Phi_c$ (or $\Phi_c^\mathbb{C}$) ensures the system sits at the threshold where the
boundary-bulk encoding is exact rather than approximate.

This suggests the following conjecture, which the grammar's precondition principle motivates
but does not yet prove.

\begin{mdframed}[style=defstyle]
\begin{questioninner}[Imscriptive involution meta-conjecture]
\label{conj:meta}
Let $(X, \tau)$ be a system where $\tau: X \to X$ is a boundary involution and the structural
type of $X$ encodes as $(\Phi_c,\ P_{\pm}^{\text{sym}},\ \Omega_{Z/Z_2})$ in the grammar.
Then:
\begin{enumerate}
  \item[(i)] There exists a distinguished generator $g \in X$ satisfying $\tau(g) = g^{-1}$
    exactly (membership in the $(-1)$-eigenspace is non-empty).
  \item[(ii)] $g$ is unique up to topological equivalence (the $(-1)$-eigenspace has rank 1
    over the appropriate coefficient ring).
  \item[(iii)] $g$ admits both an algebraic characterization (minimal polynomial / reciprocal
    relation) and an analytic characterization (L-function / functional equation / partition
    function), and these characterizations agree.
\end{enumerate}
\end{questioninner}
\end{mdframed}

\medskip

The arithmetic case is the Stark conjecture (H1 + H2a + H2b). The topological code case is
essentially the stabiliser formalism. The mirror symmetry case is the mirror conjecture. In
each domain the conjecture is either proven or a central open problem. The grammar claims they
are all instances of the same structural fact.

\subsection{The Ouroboros conjecture}

The following conjecture is more speculative. It connects the truth of Stark's conjecture to a
global completeness property of the arithmetic crystal.

\begin{mdframed}[style=defstyle]
\begin{questioninner}[Arithmetic Ouroboros]
\label{conj:ouroboros}
If Stark's conjecture holds for all ray class fields $K_d$ over real quadratic bases, then:
\begin{enumerate}
  \item[(i)] Every Stark generator $\epsd$ is a primitive element of the $O_\infty$ tier of
    the arithmetic crystal of $K_d/\bQ$.
  \item[(ii)] The minimal polynomial of $\epsd$ is the unique self-reciprocal generator of
    the $O_\infty$ Galois orbit of $K_d$.
  \item[(iii)] There exists an injective map $\iota$ from the $O_\infty$ tier of the
    arithmetic crystal to the $O_\infty$ tier of the grammar's Crystal of Types, which
    is surjective when restricted to the sub-tier of systems with $D = D_\triangle$,
    $T = T_\text{network}$, and $\Omega = \Omega_Z$.
\end{enumerate}
In other words: the Stark conjecture, if true, would imply that arithmetic is imscriptively
closed at the $O_\infty$ level — the $O_\infty$ arithmetic crystal embeds canonically into the
grammar.
\end{questioninner}
\end{mdframed}

\begin{remark}
The Ouroboros structure in (iii) is the arithmetic analog of the imscriptive monitor condition:
the boundary of the arithmetic system (the L-function, encoding analytic data) is exactly dual
to its bulk (the unit group, encoding algebraic data), with no remainder and no information
loss. This is only possible in the $O_\infty$ tier, where the Frobenius condition
$\mu \circ \delta = \mathrm{id}$ guarantees the round-trip encoding is exact. In all lower
tiers, the boundary-bulk correspondence is approximate.
\end{remark}

% ─────────────────────────────────────────────────────────────────────────────
\section{Relation to the grammar's precondition claim}
\label{sec:precondition}
% ─────────────────────────────────────────────────────────────────────────────

Section~\ref{sec:preamble} stated the strong claim: the grammar preconditions reality for
mathematics. The SIC analysis sharpens this.

The grammar's $O_\infty$ conditions ($\Phi_c^\mathbb{C}$ and $P_{\pm}^{\text{sym}}$) are the
structural necessary and sufficient conditions for a measurement to be a SIC-POVM. This is not
an analogy — the Frobenius condition in the grammar \emph{is} the Frobenius condition that
categorical quantum mechanics uses to define SICs as special commutative Frobenius algebras.

The Stark unit analysis adds precision: the grammar's type space identifies the exact algebraic
locus (the $(-1)$-eigenspace of $\tauinf$ in the unit lattice of $L_d$) where the SIC
precondition must be satisfied. The number-theoretic machinery (ray class fields, L-functions,
Galois descent) is the formal trace of that locus being instantiated. Mathematics finds
SIC-POVMs because reality already has $O_\infty$ structures at those arithmetic coordinates —
and the grammar is the theory of those coordinates.

% ─────────────────────────────────────────────────────────────────────────────

\begin{thebibliography}{9}

\bibitem{synthref}
Lando Mills.
\emph{The SynthOmnicon: A 12-Primitive Structural Grammar for Self-Organizing Systems}.
Unpublished manuscript, 2026.

\bibitem{companion}
Lando Mills.
\emph{SIC-POVM Existence via Frobenius-Fixed Stark Units}.
Submitted, April 2026.

\bibitem{afmy2017}
M.~Appleby, S.~Flammia, G.~McConnell, and J.~Yard.
SICs and algebraic number theory.
\emph{Foundations of Physics}, 47(8):1042--1059, 2017.
arXiv:1701.05200.

\bibitem{kopp2018}
Gene~S. Kopp.
SIC-POVMs and the Stark conjectures.
\emph{International Mathematics Research Notices}, 2021(18):13812--13838, 2021.
arXiv:1807.05877.

\bibitem{stark1971}
H.~M. Stark.
Values of $L$-functions at $s=1$, I.
\emph{Advances in Mathematics}, 7:301--343, 1971.

\bibitem{tate1984}
J.~Tate.
Les conjectures de Stark sur les fonctions $L$ d'Artin en $s=0$.
\emph{Progress in Mathematics}, vol.~47, Birkh\"{a}user, 1984.

\end{thebibliography}

\end{document}
